% Analogías de las lecciones de Razonamiento lógico (libros 1 y 2). Salen solas después del reto de cada lección.
% --- Libro 1 ---
\analogiade{pu-cm-1}{El edificio con sótanos}{Piensa en un edificio: la planta baja es el cero, los pisos de arriba son
los positivos y los sótanos los negativos. Sumar es subir en el ascensor; restar es bajar. Y la jerarquía de operaciones
es como vestirse: nadie se pone los zapatos antes que las medias. Primero lo de adentro (los paréntesis), después lo
demás, en su orden.}
\analogiade{pu-cm-2}{La pizza y el cuchillo}{El denominador es el cuchillo: dice en cuántos pedazos iguales se cortó la
pizza. El numerador es el plato: cuántos pedazos te serviste. Por eso no puedes sumar tercios con cuartos tal como
vienen: son pedazos de tamaños distintos. Primero hay que volver a cortar para que todos queden iguales.}
\analogiade{pu-cm-3}{Tres idiomas, un mismo mensaje}{Media pizza se puede decir en tres idiomas: «un medio» (fracción),
«cero punto cinco» (decimal) y «cincuenta por ciento» (porcentaje). El mensaje no cambia; cambia el idioma. Un buen
lector de números es trilingüe: pasa de uno a otro según cuál haga más fácil la cuenta.}
\analogiade{pu-cm-4}{El papel que se dobla}{Dobla una hoja por la mitad: dos capas. Otra vez: cuatro. Otra: ocho. Cada
doblez no \emph{suma} capas, las \emph{duplica}: eso es una potencia. La raíz hace el camino de vuelta: si tres dobleces iguales dieron 8 capas, ¿por cuánto multiplicó cada doblez? Por 2, porque $2 \cdot 2 \cdot 2 = 8$. Y estimar es mirar el grosor del papel sin contar capa por capa.}
\analogiade{pu-ar-1}{La película en reversa}{Si ves una película al revés, el vaso roto se arma y el agua vuelve a la
jarra. Así se resuelven estos problemas: se parte del final («le quedaron 5») y se rebobina. Lo que en la historia se
quitó, al rebobinar se devuelve; lo que se partió a la mitad, se duplica.}
\analogiade{pu-ar-2}{La cuenta del restaurante}{Cinco amigos piden platos de precios distintos y deciden pagar por
partes iguales: lo que paga cada uno es el promedio. Nadie regala ni pierde plata; solo se \emph{nivela}. Por eso, si
conoces el promedio y cuántos eran, conoces el total de la cuenta, y con el total se resuelve casi todo.}
\analogiade{pu-ar-3}{El detective de las cifras}{Una suma con letras es una escena con huellas. El detective no adivina:
mira la pista más reveladora (la columna de la izquierda, una llevada obligada) y de ahí deduce lo demás. Cada cifra
descubierta es un sospechoso menos.}
\analogiade{pu-ar-4}{El amigo pesimista}{Para saber si algo es \emph{seguro}, pregúntale al amigo pesimista: «¿qué es
lo peor que podría pasar?». Si incluso con la peor suerte del mundo lo logras, entonces es seguro. Con buena suerte
bastan dos medias para un par; el pesimista saca primero una de cada color.}
\analogiade{pu-op-1}{La receta inventada}{Un operador nuevo es una receta que alguien inventó: «toma el primero, elévalo
al cuadrado y réstale el doble del segundo». No importa si te parece rara: se sigue al pie de la letra. Y, como en la
cocina, cambiar el orden de los ingredientes puede cambiar el plato.}
\analogiade{pu-op-2}{Las muñecas rusas}{Un operador dentro de otro es una muñeca rusa. Si conoces lo de adentro, armas
de adentro hacia afuera. Si lo de adentro es la incógnita, toca al revés: abres primero la muñeca grande y vas llegando
a la pequeña.}
\analogiade{pu-pr-1}{La receta de la limonada}{Una limonada lleva 2 vasos de jugo por cada 3 de agua. Para una fiesta
no cambias la receta: la repites. Cuatro de jugo, seis de agua; veinte de jugo, treinta de agua. La razón $2:3$ es la
receta, y el número de veces que la repites es «cuánto vale una parte».}
\analogiade{pu-pr-2}{La torta y los invitados}{Más harina, más torta: suben juntas (directa). Pero con la misma torta,
más invitados significa pedazos más pequeños: una sube y la otra baja (inversa). Antes de hacer una regla de tres,
pregúntate cuál de las dos historias estás viviendo.}
\analogiade{pu-pr-3}{Dos mangueras, una piscina}{Dos mangueras llenando la misma piscina no promedian sus tiempos:
\emph{juntan sus chorros}. Lo que se suma es cuánto llena cada una por hora. Lo mismo pasa con dos pintores, dos
impresoras o dos carros que se acercan: se suma el ritmo, no el reloj.}
\analogiade{pu-pr-4}{El estadio de cien sillas}{Imagina cualquier cantidad convertida en un estadio de 100 sillas. El
30\,\% son 30 sillas ocupadas. La pregunta clave es siempre «¿cien sillas de \emph{qué} estadio?»: la base del
porcentaje. Cambia el estadio y las mismas 30 sillas ya no significan lo mismo.}
\analogiade{pu-pr-5}{La escalera que cambia de tamaño}{Subir un 20\,\% y luego bajar un 20\,\% es subir un escalón y
bajar otro… pero el de bajada se mide desde más arriba, así que es más alto. Por eso no vuelves al mismo piso. Los
porcentajes seguidos no se suman: cada uno se calcula sobre el piso donde quedaste.}
\analogiade{pu-ec-1}{El traductor}{Plantear una ecuación es traducir del español al idioma de las matemáticas. «El
doble» se traduce $2x$; «disminuido en 5», $-5$; «es», $=$. Un buen traductor no traduce palabra por palabra sin
pensar: primero entiende la frase completa y decide quién es $x$.}
\analogiade{pu-ec-2}{La foto familiar}{Mira dos fotos de la misma familia, con diez años de diferencia. Todos
envejecieron lo mismo: la diferencia de edad entre madre e hija es idéntica en las dos. Ese dato que no cambia es el
ancla de los problemas de edades. En las mezclas, el ancla es la cantidad de «sustancia» que no se pierde.}
\analogiade{pu-ec-3}{Dos pistas, un culpable}{Una sola pista deja muchos sospechosos: $x + y = 10$ lo cumplen
infinitas parejas. Con una segunda pista independiente queda uno solo. Y una desigualdad no señala a una persona sino
a un barrio entero: todos los que viven de cierta calle en adelante.}
\analogiade{pu-nu-1}{Los bloques de construcción}{Los números primos son las fichas básicas: no se pueden desarmar.
Todo número es una construcción hecha con esas fichas, y cada número tiene \emph{una sola} lista de fichas. El 12 es
$2\cdot2\cdot3$ y nada más. Conocer las fichas de un número es conocer todos sus divisores.}
\analogiade{pu-nu-2}{Los buses y las baldosas}{Dos buses salen juntos y uno pasa cada 12 minutos y el otro cada 18:
¿cuándo vuelven a coincidir? Eso es un \emph{múltiplo} común (el MCM, mira hacia adelante). ¿Cuál es la baldosa más
grande que cubre exacto un piso de $12 \times 18$? Eso es un \emph{divisor} común (el MCD, parte hacia adentro).}
\analogiade{pu-nu-3}{El reloj que da vueltas}{Son las 10 y pasan 5 horas: no son las 15, son las 3. El reloj solo
guarda el residuo. Con los días de la semana y con la última cifra de las potencias pasa igual: todo da vueltas en un
ciclo, y lo único que importa es en qué punto de la vuelta quedas.}
\analogiade{pu-ge-1}{La puerta que gira}{Un ángulo es cuánto se abre una puerta: nada (0°), en escuadra (90°), de par
en par (180°). No importa si la puerta es grande o pequeña: la abertura es la misma. Y dos calles paralelas cruzadas
por una avenida forman en cada esquina las mismas dos aberturas, repetidas.}
\analogiade{pu-ge-2}{El atajo por el parque}{Para ir de una esquina del parque a la opuesta puedes rodearlo (dos
catetos) o cruzarlo en diagonal (la hipotenusa). El atajo siempre es más corto que rodear, pero más largo que cualquiera
de los dos lados. Pitágoras dice exactamente cuánto mide ese atajo.}
\analogiade{pu-ge-3}{El zoom de la foto}{Al hacer zoom en una foto todo crece en la misma proporción: nadie queda con
la cabeza más grande que el cuerpo. Eso es semejanza. Pero ojo: si duplicas el zoom, la foto ocupa \emph{cuatro} veces
más pantalla, porque crece a lo ancho y a lo alto a la vez.}
\analogiade{pu-ge-4}{La cerca y el pasto}{El perímetro es la cerca que rodea el terreno; el área es el pasto que hay
adentro. Con la misma cerca se pueden encerrar potreros con mucho o con poco pasto. Por eso son dos preguntas distintas,
con unidades distintas: metros para la cerca, metros cuadrados para el pasto.}
\analogiade{pu-ge-5}{La galleta y el cortador}{El área sombreada es lo que queda de la masa después de usar el cortador
de galletas. Nadie mide esas sobras directamente: se calcula la masa completa y se le resta lo que se llevó el
cortador. Figura grande menos figura quitada.}
\analogiade{pu-ge-6}{La resma de papel}{Una hoja casi no tiene grosor, pero quinientas apiladas forman un bloque. El
volumen de un prisma es eso: el área de una hoja (la base) por cuántas hojas hay (la altura). Un cono es una pila de
hojas cada vez más pequeñas, y por eso llena solo la tercera parte.}
% --- Libro 2 ---
\analogiade{pu-or-1}{La foto del equipo}{Organizar una fila es acomodar a la gente para la foto del equipo: primero
pones a los que tienen puesto fijo («el capitán en el centro»), luego a los que van pegados («estos dos juntos») y al
final los demás van llenando huecos. Si alguien no cabe, no es su puesto.}
\analogiade{pu-or-2}{El carrusel}{En un carrusel, si todos los caballos giran, nadie cambió de vecino. En una mesa
redonda pasa igual: no hay «primer puesto». Por eso clavas a una persona en su silla, como si detuvieras el carrusel,
y ubicas a las demás respecto a ella.}
\analogiade{pu-or-3}{El sudoku}{Una tabla de doble entrada funciona como un sudoku: cada fila y cada columna aceptan
un solo acierto. Cada pista tacha casillas, y cuando en una fila queda una sola casilla libre, no hay nada que
pensar: esa es.}
\analogiade{pu-or-4}{Las reglas del parqueadero}{«Si llueve, el parqueadero cubre el carro» no dice qué pasa cuando no
llueve. Las condiciones son reglas de un reglamento: solo obligan en el caso que mencionan. Para saber qué \emph{tiene}
que pasar, hay que revisar todos los casos permitidos; para saber qué \emph{puede} pasar, basta encontrar uno.}
\analogiade{pu-lo-1}{Los interruptores}{«Y» son dos interruptores en fila: la luz prende solo si los dos están
cerrados. «O» son dos caminos distintos hacia el bombillo: basta uno. «No» es el interruptor que hace lo contrario de
lo que le dicen. Con esas tres piezas se arma cualquier razonamiento.}
\analogiade{pu-lo-2}{El abogado contrario}{Negar una frase es hacer de abogado contrario: ¿qué es lo \emph{mínimo}
que necesito para demostrar que mientes? Si dices «todos llegaron», me basta \emph{uno} que no llegó. Si dices «estudia
y trabaja», me basta que falle una de las dos.}
\analogiade{pu-lo-3}{La promesa}{«Si ganas el año, te regalo la bicicleta» es una promesa. Solo se rompe de una
manera: ganaste el año y no hubo bicicleta. Si no ganaste, la promesa no se rompió, haya regalo o no. Y recibir una
bicicleta no prueba que ganaste el año: pudo ser tu cumpleaños.}
\analogiade{pu-lo-4}{Las cajas dentro de cajas}{«Todos los paisas son colombianos» es una caja pequeña dentro de una
grande. Quien está en la pequeña está en la grande; pero estar en la grande no te mete en la pequeña. Dibujar las
cajas es más seguro que confiar en cómo suena la frase.}
\analogiade{pu-lo-5}{El juego de «supongamos»}{Con mentirosos y veraces se juega a «supongamos que este dice la
verdad» y se sigue el hilo. Si el hilo se enreda (alguien resulta veraz y mentiroso a la vez), la suposición era
falsa. Es como probar una llave: si no abre, pruebas la otra.}
\analogiade{pu-cj-1}{Los grupos de WhatsApp}{Tus contactos son el universo. El grupo de la familia y el del colegio
son dos conjuntos. Los primos que estudian contigo están en los dos: la intersección. Todos los que están en al menos
uno: la unión. Y los que no están en ninguno son el complemento.}
\analogiade{pu-cj-2}{El censo sin contar dos veces}{Si sumas los que juegan fútbol con los que juegan baloncesto, a
quienes juegan ambos los contaste dos veces. Por eso el diagrama se llena desde el centro: primero los de ambos, y a
cada grupo se le resta lo que ya quedó puesto.}
\analogiade{pu-co-1}{El menú del día}{Un almuerzo tiene sopa, seco y jugo. Con 2 sopas, 3 secos y 2 jugos no hay
$2 + 3 + 2$ almuerzos sino $2 \cdot 3 \cdot 2$: cada sopa se combina con cada seco y con cada jugo. Cuando eliges una
cosa \emph{y} otra, multiplicas; cuando eliges una cosa \emph{o} otra, sumas.}
\analogiade{pu-co-2}{El podio}{En una carrera no da lo mismo quién queda de primero y quién de tercero: oro, plata y
bronce son puestos distintos. Cuando el orden importa, cada reordenamiento cuenta como un resultado nuevo. Eso es una
permutación.}
\analogiade{pu-co-3}{La ensalada de frutas}{En una ensalada de mango, piña y fresa da igual cuál fruta echaste
primero: es la misma ensalada. Cuando solo importa \emph{quiénes} quedaron elegidos y no en qué orden, es una
combinación. Por eso hay menos combinaciones que permutaciones: muchas listas ordenadas son la misma ensalada.}
\analogiade{pu-co-4}{Contar las sillas vacías}{Para saber cuántos fueron al cine, a veces es más fácil contar las
sillas vacías y restarlas del total. Con «al menos uno» pasa lo mismo: en vez de contar todos los casos buenos, cuenta
el único caso malo («ninguno») y réstalo.}
\analogiade{pu-co-5}{Las calles de la ciudad}{Ir de la casa al colegio por una ciudad en cuadrícula, sin devolverse,
es decidir en cada esquina «derecha» o «arriba». Cada ruta es una palabra hecha con esas dos letras. Contar rutas es
contar palabras, y eso ya lo sabes hacer.}
\analogiade{pu-pb-1}{La rifa}{En una rifa de 100 boletas, si compras 3, tu probabilidad es 3 de 100. Así de simple:
boletas tuyas sobre boletas totales. La única condición es que todas las boletas tengan la misma oportunidad; por eso
con dos dados se cuentan las 36 parejas, no las 11 sumas.}
\analogiade{pu-pb-2}{La bolsa de dulces}{Si sacas un dulce y lo devuelves, la bolsa queda igual: la segunda sacada no
se entera de la primera (independientes). Si te lo comes, la bolsa cambió y las cuentas también. Y para «al menos
uno», acuérdate de las sillas vacías: uno menos la probabilidad de «ninguno».}
\analogiade{pu-se-1}{El ritmo de una canción}{Una sucesión tiene ritmo. Unas avanzan a paso parejo (suman siempre lo
mismo), otras se aceleran (multiplican), otras alternan dos compases. Antes de adivinar el número que sigue, escucha
el ritmo: resta vecinos, divide vecinos, y mira si hay dos melodías entrelazadas.}
\analogiade{pu-se-2}{El alfabeto es una regla}{Las letras son números disfrazados: A es 1, B es 2, C es 3… Quítales
el disfraz y una sucesión de letras se vuelve una de números. Con las figuras, el truco es mirar una sola cosa a la
vez: primero cómo gira, después cuántas hay, después el color.}
\analogiade{pu-se-3}{El semáforo}{Verde, amarillo, rojo, y vuelve a empezar. Para saber de qué color estará dentro de
mil cambios no hay que contarlos todos: basta quitar las vueltas completas y mirar lo que sobra. El residuo es el
mapa de cualquier cosa que se repite.}
\analogiade{pu-cd-1}{El diccionario secreto}{Un código es un diccionario secreto: cada letra tiene su traducción. Si
descubres que la M se volvió N y la E se volvió F, ya tienes la regla del diccionario (correr una letra) y puedes
traducir cualquier palabra, la hayas visto o no.}
\analogiade{pu-cd-2}{Aprender un idioma viendo subtítulos}{Si en dos escenas aparece la palabra «sapi» y en las dos el
subtítulo dice «grande», ya aprendiste una palabra. Así se descifra un idioma inventado: buscando lo que se repite a
ambos lados. Lo que cambia junto, significa lo mismo.}
\analogiade{pu-es-1}{El GPS interior}{Seguir instrucciones en una rejilla es ser tu propio GPS: en cada paso debes
saber dónde estás y hacia dónde miras. «Gira a la derecha» no te mueve del sitio, solo cambia hacia dónde apunta tu
nariz. Y la derecha siempre es la \emph{tuya}, no la de quien mira el mapa.}
\analogiade{pu-es-2}{El edificio visto desde un dron}{Desde un dron, un edificio de 20 pisos y uno de 2 se ven como
dos rectángulos iguales. Por eso la vista de arriba necesita números que digan la altura de cada torre. Contar cubos
es contar apartamentos: torre por torre, incluidos los pisos que el dibujo tapa.}
\analogiade{pu-es-3}{La caja de pizza desarmada}{Una caja de cartón desarmada es un desarrollo: caras planas unidas
por dobleces. Al armarla, las caras que eran vecinas en el cartón siguen siendo vecinas, y las que estaban «a una de
distancia» terminan frente a frente.}
\analogiade{pu-es-4}{La mancha de tinta}{Dobla una hoja con una gota de tinta fresca y ábrela: la mancha aparece
repetida al otro lado del doblez, como en un espejo. Cada doblez es un espejo. Para saber cómo queda la hoja,
desdobla de a un espejo, empezando por el último.}
\analogiade{pu-gr-1}{El tablero del carro}{Una gráfica es el tablero de un carro: muestra mucho en poco espacio, pero
solo si sabes qué mide cada aguja. Antes de leer el número, mira el nombre del reloj y sus unidades. Confundir el
velocímetro con el medidor de gasolina es el error más común, también en el examen.}
\analogiade{pu-gr-2}{Tres voceros del grupo}{Para hablar por un grupo hay tres voceros. El \emph{promedio} reparte
todo por igual. La \emph{mediana} es el que queda en la mitad de la fila ordenada. La \emph{moda} es lo que más se
repite. Si al grupo entra un millonario, el promedio se dispara y la mediana casi ni se entera.}
\analogiade{pu-gr-3}{El testigo que solo dice lo que vio}{Una gráfica es un testigo honesto pero limitado: dice lo que
midió y nada más. Si muestra que subieron las ventas de helado y el calor, eso vio. Que una cosa \emph{cause} la otra,
o que pase igual en otra ciudad, ya no lo dice el testigo: lo estarías inventando tú.}
